% 分野別類題: 熱伝導方程式（変数分離法）
% u_t = k u_xx （有限区間、ディリクレ境界条件2問＋ノイマン境界条件1問）
% コンパイル: TEXINPUTS="../../../00_common:$TEXINPUTS" latexmk -xelatex problems.tex
\documentclass[a4paper,11pt]{article}
\usepackage{preamble}

\begin{document}

\kamokumei{類題演習 --- 熱伝導方程式（変数分離法）}

\shijibun{次の熱伝導方程式の初期値・境界値問題を、変数分離法により解け。導出過程を示すこと。}

\shomon{問1.}{次の初期値・境界値問題を解け。}
\[
u_t = k\,u_{xx} \qquad (0 < x < L,\ t > 0)
\]
\[
u(0,t) = 0,\quad u(L,t) = 0 \qquad (t > 0)
\]
\[
u(x,0) = \sin\frac{\pi x}{L} \qquad (0 \le x \le L)
\]

\shomon{問2.}{次の初期値・境界値問題を解け。}
\[
u_t = k\,u_{xx} \qquad (0 < x < L,\ t > 0)
\]
\[
u(0,t) = 0,\quad u(L,t) = 0 \qquad (t > 0)
\]
\[
u(x,0) = 3\sin\frac{\pi x}{L} - 2\sin\frac{3\pi x}{L} \qquad (0 \le x \le L)
\]

\shomon{問3.}{両端が断熱された棒における次の初期値・境界値問題を解け。}
\[
u_t = k\,u_{xx} \qquad (0 < x < L,\ t > 0)
\]
\[
u_x(0,t) = 0,\quad u_x(L,t) = 0 \qquad (t > 0)
\]
\[
u(x,0) = u_0 + \cos\frac{\pi x}{L} \qquad (0 \le x \le L,\ u_0 \text{ は定数})
\]

\end{document}
